\section{Experiments}
\label{sec:experiments}

The evaluation is organized around three questions. \textbf{Q1:} Do correction
intervals and successful nominal segments provide reliable and complementary
supervision? \textbf{Q2:} Does the learned filter preserve adequate human
motion while improving difficult states and unseen configurations?
\textbf{Q3:} Does iterative updating increase valid demonstrations per unit of
operator time, and do those demonstrations improve a downstream policy?
Unless stated otherwise, comparisons use the same robot, control interface,
active operator time, reset count, configuration quota, safety limits, and
episode-level data split.

\subsection{Setup and Evaluation Protocol}
\label{sec:exp_setup}

Each episode contains synchronized RGB observations, raw, filtered, and
executed commands, robot state, audit events, correction intervals, terminal
status, and configuration metadata. Tasks start with a grasped object and a
visually observable target, and include local approach, alignment, correction,
or short-horizon execution. Success tolerances, velocity limits, and safety
termination conditions are fixed before collection. Participants receive the
same familiarization period and use every compared method on matched
configurations in balanced crossover order. The primary collection-efficiency
analysis is paired within novice operators; experienced operators provide a
secondary experience-stratified analysis.

\begin{table}[!t]
  \centering
  \caption{Experimental protocol.}
  \label{tab:protocol}
  \footnotesize
  \setlength{\tabcolsep}{4pt}
  \begin{tabularx}{\columnwidth}{@{}>{\raggedright\arraybackslash}p{0.43\columnwidth}Y@{}}
    \toprule
    Item & Configuration \\
    \midrule
    Robot and end-effector & \\
    Teleoperation interface & \\
    Cameras and placement & \\
    Image resolution / rate & \\
    Control mode / rate & \\
    Synchronization error & \\
    Visual encoder & SigLIP2: \\
    Offline auditor & Qwen-VL: \\
    History $L$; gain limits $\alpha_{\max},r_\alpha$ & \\
    Tasks / configurations & \\
    Episodes per configuration & \\
    Operators (novice / experienced) & \\
    Familiarization / crossover order & \\
    Train / validation / test episodes & \\
    $H_{\mathrm{ref}}$ / round-audit episodes & \\
    Operator-time / reset / safety budgets & \\
    \bottomrule
  \end{tabularx}
\end{table}

All splits are made at the complete-episode level, so adjacent windows from a
single trajectory never cross train, validation, and test sets. The cold-start
round uses manually confirmed action data. Later rounds apply the same
admission rule $\Gamma$ to update $\Aaction$; failed, near-miss, conflicting,
and anomalous episodes remain in $\Aaudit$. Metrics are aggregated within an
episode and then across operators and configurations. We use
\begin{align}
  E_{\mathrm{nom}}&=\frac{1}{|\mathcal{N}|}\sum_{t\in\mathcal{N}}
    \|u^{\mathrm{out}}_t-u^{\mathrm{raw}}_t\|_1,\qquad
  G_{\mathrm{corr}}=\frac{1}{|\mathcal{C}|}\sum_{t\in\mathcal{C}}\alpha_t,\\
  \Delta G&=G_{\mathrm{corr}}-G_{\mathrm{nom}},\qquad
  G_{\mathrm{nom}}=\frac{1}{|\mathcal{N}|}\sum_{t\in\mathcal{N}}\alpha_t,
  \label{eq:exp_metrics}
\end{align}
where $\mathcal{C}=\{t:m_t=1\}$ and $\mathcal{N}=\{t:m_t=0\}$. We also
report action prediction error, gain AUPRC, gain total variation, jerk, latency,
clipping, bypass, and safety events. A threshold $\eta$ is selected on the
validation set and frozen for interval IoU and boundary-delay statistics.

\subsection{Auditable Supervision}
\label{sec:exp_supervision}

Q1 tests both the audit signal and the information supplied by its two episode
parts. A second annotator independently labels correction intervals and
terminal outcomes in the frozen reference set $H_{\mathrm{ref}}$. We report
interval IoU, active-frame F1, boundary error, success/failure agreement,
admission agreement, and review time for manual labels, VLM proposals with
review, and unreviewed proposals. Table~\ref{tab:q1} retains the primary audit
and assistance outcomes; boundary and admission errors are reported alongside
the corresponding intervals in the text. We then train the filter with full,
nominal-only, correction-only, and reviewed-VLM supervision. Action prediction
error and gain AUPRC are secondary diagnostics.

\begin{table}[!t]
  \centering
  \caption{Audit reliability and complementary-supervision ablations.}
  \label{tab:q1}
  \footnotesize
  \setlength{\tabcolsep}{3pt}
  \textbf{(a) Audit reliability and effort}\par\smallskip
  \begin{tabularx}{\columnwidth}{@{}Yccc@{}}
    \toprule
    Audit procedure & IoU $\uparrow$ & F1 $\uparrow$ & Review time $\downarrow$ \\
    \midrule
    Independent manual relabel & & & \\
    VLM proposals + review & & & \\
    VLM proposals & & & \\
    \bottomrule
  \end{tabularx}
  \smallskip
  \textbf{(b) Supervision ablation}\par\smallskip
  \begin{tabularx}{\columnwidth}{@{}Yccc@{}}
    \toprule
    Supervision & $E_{\mathrm{nom}}\downarrow$ & $\Delta G\uparrow$ & $S_{\mathrm{exec}}\uparrow$ \\
    \midrule
    Full complementary supervision & & & \\
    Nominal-only & & & \\
    Correction-only & & & \\
    Reviewed VLM supervision & & & \\
    \bottomrule
  \end{tabularx}
\end{table}

The comparison is intended to show whether nominal successes constrain
unnecessary offsets while correction intervals improve local recovery, rather
than merely increasing the number of training windows.

\subsection{Continuous Assistance and Generalization}
\label{sec:exp_filter}

Q2 compares raw teleoperation, a fixed temporal filter, an action-only
trajectory predictor, and the complete filter. All methods use the same safety
projection and data budget. Results are split between successful nominal
segments and episodes containing correction intervals. We then remove visual
context, merge the action and gain branches, or remove the gain-rate
constraint. The frozen model is then evaluated on unseen object poses and
configurations. Table~\ref{tab:q2} denotes gain total variation by
$V_\alpha$ and execution success by $S_{\mathrm{exec}}$.

\begin{table}[!t]
  \centering
  \caption{Continuous assistance, architecture ablations, and held-out
  generalization.}
  \label{tab:q2}
  \footnotesize
  \setlength{\tabcolsep}{2.5pt}
  \begin{tabularx}{\columnwidth}{@{}Yccccc@{}}
    \toprule
    Method / setting & $E_{\mathrm{nom}}\downarrow$ & $\Delta G\uparrow$ & $V_\alpha\downarrow$ & Jerk $\downarrow$ & $S_{\mathrm{exec}}\uparrow$ \\
    \midrule
    \multicolumn{6}{@{}l}{\itshape Baselines} \\
    Raw teleoperation & & & & & \\
    Fixed temporal filter & & & & & \\
    Action-only trajectory model & & & & & \\
    \method{} (ours) & & & & & \\
    \addlinespace[2pt]
    \multicolumn{6}{@{}l}{\itshape Architecture ablations} \\
    Without visual context & & & & & \\
    Shared action/gain branch & & & & & \\
    Without gain-rate limit & & & & & \\
    \addlinespace[2pt]
    \multicolumn{6}{@{}l}{\itshape Generalization} \\
    In-distribution & & & & & \\
    Unseen configuration & & & & & \\
    \bottomrule
    \end{tabularx}
\end{table}

The primary interpretation is a paired trade-off: $E_{\mathrm{nom}}$ measures
how closely the filter preserves adequate operator input, while $\Delta G$,
$V_\alpha$, jerk, and $S_{\mathrm{exec}}$ characterize assistance around
difficult states. We report means with standard deviations or 95\% confidence
intervals; the final caption specifies the repeated-measures test and unit of
analysis.

\subsection{Fixed-Budget Collection and Downstream Utility}
\label{sec:exp_flywheel}

Q3 compares a static filter, single-round training, and iterative updates under
matched operator active time, reset count, configuration quota, and safety
budget. Every round uses the frozen reference set and admission rule. A
mutually exclusive $H^{(r)}_{\mathrm{audit}}$ is manually reviewed to estimate
admission precision, recall, and conflict rate on newly collected episodes.
The primary quantity is valid-demonstration yield,
\begin{equation}
  \rho_r=\frac{|\Aaction^{(r)}|}
  {\operatorname{operator\ minutes}^{(r)}}.
  \label{eq:yield}
\end{equation}
We train one fixed downstream learner, ACT or Diffusion Policy, with identical hyperparameters on the
admitted data from each round \citep{zhaoLearningFineGrainedBimanual2023,chiDiffusionPolicyVisuomotor2024}.
The downstream split is disjoint from filter training, and downstream results
do not affect data admission.

\begin{table*}[!t]
  \centering
  \caption{Fixed-budget flywheel and downstream utility.}
  \label{tab:q3}
  \footnotesize
  \setlength{\tabcolsep}{3pt}
  \begin{minipage}[t]{0.47\textwidth}
    \centering
    \textbf{(a) Collection yield and audit stability}\par\smallskip
    \begin{tabularx}{\linewidth}{@{}Ycccc@{}}
      \toprule
      Procedure & $\rho_0\uparrow$ & $\rho_1\uparrow$ & $\rho_2\uparrow$ & Audit P/R $\uparrow$ \\
      \midrule
      Static filter & & & & \\
      Single-round training & & & -- & \\
      Iterative flywheel & & & & \\
      \bottomrule
    \end{tabularx}
  \end{minipage}\hfill
  \begin{minipage}[t]{0.47\textwidth}
    \centering
    \textbf{(b) Downstream utility}\par\smallskip
    \begin{tabularx}{\linewidth}{@{}Yccc@{}}
      \toprule
      Data source & Episodes & $S_{\mathrm{test}}\uparrow$ & $N_{\mathrm{target}}\downarrow$ \\
      \midrule
      Raw teleoperation & & & \\
      Static-filter data & & & \\
      Iteration 1 admitted data & & & \\
      Iteration 2 admitted data & & & \\
      \bottomrule
    \end{tabularx}
  \end{minipage}
\end{table*}

Here $S_{\mathrm{test}}$ is held-out policy success and
$N_{\mathrm{target}}$ is the number of admitted episodes needed to reach the
predeclared target performance. Safety events are reported with operator-level
counts and exposure time rather than compressed into the utility table.
Failure, near-miss, recovery, and conflict types are reported separately from
$\Aaudit$ to determine whether iteration changes the composition of rejected
episodes as well as the admitted count. The main claim is supported only when
$\rho_r$ rises without a corresponding drop in audit precision or downstream
utility.
